\documentclass[11pt,a4paper]{article}
\usepackage[T1]{fontenc}
\usepackage[utf8]{inputenc}
\usepackage[italian]{babel}
\usepackage{amsmath,amssymb}
\usepackage{geometry}
\geometry{margin=1.15in}
\usepackage{booktabs}
\usepackage[expansion=false]{microtype}
\usepackage[colorlinks=true,linkcolor=black,urlcolor=blue!50!black]{hyperref}
\usepackage{xcolor}

\title{Stesse scelte, ragionamenti diversi\\[2pt] \large Nota di orientamento sul progetto}
\author{Francesco Pinna}
\date{16 ottobre 2026}

\begin{document}
\maketitle

\section{La domanda}

Un'offerta eccessiva per un bene di qualit\`a incerta ammette due spiegazioni. La prima \`e una valutazione troppo alta del bene. La seconda \`e la trascuratezza della selezione operata dal venditore, che accetta di vendere pi\`u volentieri quando il bene vale poco. L'offerta finale da sola non separa le due spiegazioni, perch\'e possono produrre lo stesso numero, mentre implicano credenze intermedie diverse e reazioni diverse all'informazione nuova. La struttura della decisione \`e quella del mercato dei limoni di Akerlof (1970) e dell'acquisto di un'impresa di Samuelson e Bazerman (1985), con il venditore selettivo nel ruolo della parte informata, e Charness e Levin (2009) mostrano che il fenomeno sopravvive anche contro una controparte programmata, che \`e la semplificazione adottata qui. L'antecedente pi\`u vicino nell'ambiente del corso \`e l'esperimento acquiring-a-company con messaggi sul valore di Di Cagno, G\"uth, Lohse, Marazzi e Spadoni (2022).

Il progetto fa due domande. La prima chiede se il divario fra il prezzo di riserva per un venditore sempre disponibile e quello per un venditore selettivo sia pi\`u piccolo del divario implicato dal condizionamento corretto, tenuto conto della valutazione individuale del bene. La seconda chiede se raccogliere decisioni intermedie cambi il prezzo finale, cos\`i che la tecnologia di misura viene studiata come un trattamento e non assunta innocua. Entrambe vengono dal preprint di ancoraggio, Pinna (2026), che nelle sezioni 10.3 e 12.4 e nelle conclusioni chiede esattamente un esperimento con tempi dei report randomizzati, un braccio passivo e report ripetuti a informazione invariata dai quali stimare il rumore di risposta.

\section{Le quattro idee teoriche che servono}

\paragraph{Il posterior come martingala.}
Indichiamo con $\pi$ la probabilit\`a corrente che il bene valga 90, un numero fra zero e uno, aggiornata con la regola di Bayes dai certificati annunciati. Il prezzo di riserva di chi \`e neutrale al rischio e condiziona correttamente, per il venditore sempre disponibile, \`e
\[
U(\pi)=30+60\pi.
\]
Poich\'e $\pi$ \`e una media condizionata, le sue revisioni hanno media zero date le informazioni gi\`a viste, e lo stesso vale per $U$. Questa propriet\`a di martingala \`e il punto fermo del disegno, perch\'e d\`a una scala osservabile alle revisioni bayesiane e un benchmark che non dipende dal modello comportamentale.

\paragraph{La contrazione cursed.}
La seconda idea \`e il cursed equilibrium di Eyster e Rabin (2005), nella forma usata dal preprint. Si prende la tabella congiunta $p$ delle probabilit\`a di valore e disponibilit\`a, si prende il prodotto delle sue marginali, e per un parametro $\chi$ fra zero e uno si definisce la tabella percepita
\[
\Gamma_{\chi}(p)=(1-\chi)\,p+\chi\,p_{V}\otimes p_{J}.
\]
Le marginali restano quelle vere e l'unica cosa che cambia \`e l'associazione fra valore e disponibilit\`a, che viene contratta. A $\chi=0$ c'\`e il condizionamento corretto, a $\chi=1$ la disponibilit\`a viene trattata come indipendente dal valore.

\paragraph{L'identit\`a in prezzi.}
La terza idea traduce la contrazione in prezzi, ed \`e tutta l'algebra che serve. Indichiamo con $a$ e $b$ le probabilit\`a che il venditore selettivo sia disponibile quando il bene vale 90 e quando vale 30. Il prezzo condizionato correttamente \`e
\[
C_{a,b}(\pi)=30+60\,\frac{a\pi}{a\pi+b(1-\pi)}.
\]
Sotto la trasformazione, il valore atteso del bene dato che il venditore \`e disponibile diventa una media pesata fra valore condizionato e valore incondizionato, quindi il prezzo selettivo \`e esattamente
\[
S(\pi,\chi)=C_{a,b}(\pi)+\chi\,\bigl[U(\pi)-C_{a,b}(\pi)\bigr],
\]
una retta fra i due benchmark. Chi trascura del tutto la selezione paga il venditore selettivo come quello sempre disponibile. Se il primo prezzo misura il posterior soggettivo attraverso $\widetilde{\pi}=(U-30)/60$, il coefficiente ricavabile dalla coppia di prezzi, che il fascicolo chiama contrazione effettiva e indica con $\psi$, \`e
\[
\psi=\frac{S-C_{a,b}(\widetilde{\pi})}{\,U-C_{a,b}(\widetilde{\pi})\,},
\]
e nel benchmark statico con posterior corretto coincide con $\chi$.

\paragraph{Misurare senza fidarsi della misura.}
La quarta idea giustifica i tre protocolli. Ci\`o che un esperimento osserva \`e una proiezione del percorso di ragionamento, e il preprint dimostra che l'informazione contenuta nell'osservazione \`e la media condizionata di quella contenuta nel percorso, cos\`i che due ragionamenti diversi possono diventare indistinguibili nella sola scelta finale. Chiedere di pi\`u, per\`o, pu\`o disturbare. Un report intermedio si scrive come il valore latente, pi\`u un eventuale effetto sistematico del protocollo di rilevazione, pi\`u un rumore transitorio. I report ripetuti a informazione ferma identificano la varianza del rumore, mentre l'effetto sistematico non si assorbe nel rumore e va trattato come un trattamento randomizzato. Per questo il disegno randomizza i protocolli, tiene un braccio senza misure intermedie e usa la calibrazione a valore certo come stima del rumore. Il problema di incentivi delle scoring rule documentato da Danz, Vesterlund e Wilson (2022) \`e la ragione per cui si chiedono prezzi e mai probabilit\`a.

\section{La struttura dell'esperimento}

In ogni prova del modulo principale il bene vale 30 oppure 90 punti, il prior annunciato sul valore alto \`e 0,4 oppure 0,6, e arrivano in ordine casuale due certificati condizionalmente indipendenti dato il valore, corretto il debole nel 60\% dei casi e il forte nel 70\%. I due venditori sono programmati e la regola \`e dichiarata, quindi non c'\`e inganno e il benchmark \`e calcolabile. Il selettivo forte \`e disponibile l'80\% delle volte se il bene vale 30 e il 20\% se vale 90, il debole il 65\% e il 35\%, il placebo il 50\% in entrambi i casi. Le sedici prove comprendono otto prove dipendenti, prodotto di due prior, due intensit\`a e due ripetizioni, quattro placebo indipendenti e quattro prove di rivelazione, nelle quali la disponibilit\`a estratta viene mostrata prima della lista finale. Per ciascun venditore si indica fino a quale prezzo si comprerebbe, su una lista a passo uno di 71 righe, nel formato 20-90 in met\`a delle prove e 30-100 nell'altra met\`a, cos\`i che il centro numerico cambi e l'ancoraggio sia verificabile. Niente \`e preselezionato, e se una lista viene pagata il computer estrae un prezzo e applica a quel prezzo la decisione dichiarata, quindi la soglia indicata non \`e mai il prezzo pagato in automatico.

\begin{table}[ht]
\centering
\begin{tabular}{lccc}
\toprule
Posterior 0,4 & Prezzo razionale & Contrazione 0,25 & Contrazione 1\\
\midrule
Sempre disponibile & 54,00 & 54,00 & 54,00\\
Selezione forte & 38,57 & 42,43 & 54,00\\
Selezione debole & 45,85 & 47,89 & 54,00\\
\bottomrule
\end{tabular}
\caption{La predizione al posterior 0,4. Il prezzo cursed \`e sulla retta fra i due benchmark.}
\end{table}

I tre protocolli differiscono solo in quando si misura. F raccoglie le sole decisioni finali, S aggiunge una rilevazione intermedia per prova, scelta fra due momenti e bilanciata su prior e intensit\`a, R raccoglie entrambe le rilevazioni. L'assegnazione \`e alla diade, bloccata per sessione, i tempi delle finestre sono fissi e uguali, una pausa di trenta secondi con un nuovo certificato separa l'ultima misura dalla finestra finale, e la probabilit\`a che una decisione venga pagata \`e la stessa nei tre protocolli, perch\'e uno slot inattivo conserva una lotteria di riserva al 50\%.

Le due ipotesi primarie sono queste. La prima \`e una contrazione effettiva media positiva nelle otto prove dipendenti del protocollo S, stimata con una likelihood che rispetta gli intervalli delle liste e il modello di misura dichiarato, con test unilaterale al livello 0,025, perch\'e dividere un prezzo rumoroso per un altro prezzo rumoroso distorcerebbe il rapporto. La seconda \`e l'equivalenza fra F e R del prezzo finale per il venditore selettivo entro tre punti, il 5\% della distanza fra i valori, verificata con la procedura dei due test unilaterali al livello 0,025 per lato, cos\`i che la famiglia delle due primarie \`e controllata al 5\%. Tutto il resto, compreso il confronto di S con gli altri due protocolli, sta nelle analisi secondarie con un controllo dichiarato del false discovery rate.

Completano il disegno otto prove a valore certo, che stimano rumore e ancoraggio dello strumento a informazione ferma, sedici urne, le prove con il computer collocate dopo i moduli principali, e un modulo esplorativo sulla profondit\`a di ragionamento con il partner reale, tre giochi selezionati e sei prove a ruoli scambiati. Il pagamento estrae una decisione per persona e il premio \`e una lotteria binaria, cos\`i i prezzi di riserva non ereditano le preferenze per il rischio. La base di pianificazione \`e di 300 persone in 150 diadi, cinquanta per protocollo, in dieci sessioni da trenta, con un pilota separato di 36 che non entra nell'analisi e con una regola dichiarata prima dei dati, poich\'e l'80\% di potenza per l'equivalenza a 300 richiede dal pilota una deviazione standard di circa 5,92 o inferiore. Una sessione dura da 85 a 95 minuti. L'esperimento \`e progettato e documentato per la presentazione, e non \`e stato eseguito.

\section{Dove trovare il resto}

Il fascicolo contiene il protocollo completo, le istruzioni ai partecipanti, il piano di analisi e i due programmi di verifica che riproducono ogni numero di questa nota, compresa la tabella qui sopra. Il documento in inglese e le slide coprono la presentazione, con una sezione di riferimenti d'ingresso pensata per le docenti. Il preprint di ancoraggio \`e pubblico, Pinna (2026), SSRN 7251718, e la nota non ne modifica alcun risultato.

\subsection*{Riferimenti}

\begin{list}{}{\setlength{\itemsep}{2pt}\setlength{\leftmargin}{1.5em}}
\item Akerlof, G. A. (1970). The Market for ``Lemons'': Quality Uncertainty and the Market Mechanism. \emph{Quarterly Journal of Economics}, 84, 488-500.
\item Charness, G., e Levin, D. (2009). The Origin of the Winner's Curse: A Laboratory Study. \emph{American Economic Journal: Microeconomics}, 1, 207-236.
\item Danz, D., Vesterlund, L., e Wilson, A. J. (2022). Belief Elicitation and Behavioral Incentive Compatibility. \emph{American Economic Review}, 112, 2851-2883.
\item Di Cagno, D., G\"uth, W., Lohse, T., Marazzi, F., e Spadoni, L. (2022). Who Cares When Value (Mis)reporting May Be Found Out? CEIS Research Paper 531. SSRN 4022439.
\item Eyster, E., e Rabin, M. (2005). Cursed Equilibrium. \emph{Econometrica}, 73, 1623-1672.
\item Pinna, F. (2026). Dynamic Identification of Behavioral Reasoning in Finite Games. SSRN Working Paper 7251718, versione dell'8 agosto 2026.
\item Samuelson, W. F., e Bazerman, M. H. (1985). The Winner's Curse in Bilateral Negotiations. In \emph{Research in Experimental Economics}, volume 3, JAI Press, 105-137.
\end{list}

\end{document}
